A feature-based and a raw-byte malware detector explain their decisions in spaces that do not meet: one in features
with no location in the file, the other in bytes with no name. We showed that recording, during feature extraction,
which bytes produce each feature gives a single provenance matrix that carries explanations both ways. Pulled back, a
feature-based detector's explanation lands on the bytes at the cost of one extraction, as faithfully as direct
problem-space explanation at the same resolution and more so at byte resolution, and lets two detectors be compared on
the same file at a resolution that windows of bytes cannot reach. Pushed forward, a raw-byte detector's explanations
acquire the extractor's feature names, which recur across files, so its reliance can be aggregated over a corpus,
confirmed, shown to be exploitable with changes the loader allows, and compared with a feature-based detector's. Every
engineered feature set is computed by a program from an input; we expect the same
matrix, recorded once per extractor, to make the explanations of detectors in other domains comparable and auditable
in the same way.
